\begin{lemma}
Let $D$ be a finite loopless digraph with no copy of $T_s$, and let $k\geq
1$.  There is a loopless graph $G$ on the same number of labelled vertices,
containing no clique of order $s$, such that
\[
 i_k(G)\leq \left(\frac{e}{k}\right)^k\operatorname{fwi}_k(D).
\]
Here the tuple count is the ordered forward-independent count of
\noderef{ForwardIndependentCount}, while $i_k$ is the unordered independent
set count of \noderef{IndependentSetCount}.
\end{lemma}
\begin{proof}
Choose a uniformly random permutation $\pi$ of $V(D)$.  Define an undirected
graph on the same vertex set by putting an edge $\{u,v\}$, for $u\ne v$,
whenever one of the two directed arcs between $u$ and $v$ points forward in
the order $\pi$; equivalently, include the arc $(u,v)$ when
$\pi(u)<\pi(v)$ and forget its orientation.  This is a loopless symmetric
graph, and it has exactly the same number of vertices as $D$.

If this graph contained a clique $u_1,\ldots,u_s$ in the exact witness sense
of \noderef{CliqueWitness}, list those vertices in increasing order of their
$\pi$-values.  For every $i<j$, the edge between $u_i$ and $u_j$ was created by
an arc $(u_i,u_j)$, so the ordered vertices would form a copy of the
transitive tournament \noderef{TransitiveTournament}, contradicting the
assumed $T_s$-freeness.

Fix an ordered forward-independent $k$-tuple of distinct vertices.  The
probability that its entries occur in increasing $\pi$-order is $1/k!$.
Whenever this happens, the entries form an independent $k$-set of the
random graph: every possible forward arc is absent by forward independence.
Conversely, every independent $k$-set has at most one increasing ordering,
so linearity of expectation gives
\[
 \mathbb E i_k(G)\leq \frac{\operatorname{fwi}_k(D)}{k!}.
\]
Repeated entries are counted by the ordered convention of
\noderef{ForwardIndependentCount} but can never be an increasing ordering of
a set counted by \noderef{IndependentSetCount}, which explains the inequality.
Finally,
$k!\geq(k/e)^k$ gives $1/k!\leq(e/k)^k$.  Therefore some permutation has
the displayed bound, and the associated graph is the required $G$.
\end{proof}
